\clearpage
\SourceLeaf{328}
\pdfbookmark[1]{Bemerkungen zu den Zusätzen}{zusaetze-bemerkungen}
\begin{center}
\textbf{BEMERKUNGEN.}
\end{center}

Die Notizen [2], [4], [5], [13], [15] finden sich in einem Handbuch, [3], [9], [10], [17], [18] auf losen Blättern, [1] und [12] in \textsc{Gauss}’ Exemplaren des Berliner Astronomischen Jahrbuchs für 1805 und 1818 und [6] in \textsc{Gauss}’ Exemplar von \textsc{Vega}, log.-trig. Handbuch, Leipzig, 1797. Die Notizen [7], [8], [11], [16] sind bereits an den angegebenen Stellen veröffentlicht. Die Notiz [14] stammt aus einem Heft mit dem Titel »Abhandlungen von \textsc{Carl Friedrich Gauss}, Erstes Heft«, das Herr Professor \textsc{Jacobi} in Göttingen freundlicherweise zur Verfügung gestellt hat und das aus dem Nachlass von \textsc{Wilhelm Weber} stammt. Im Einband findet sich die Namenseintragung »\textsc{J. H. Westphal}, Danzig 1818«. Außer einigen Kopien von Aufsätzen \textsc{Gauss}’ enthält das Heft Notizen derselben Art wie die Notiz [14], die aus Vorlesungen oder anderen mündlichen Mitteilungen von \textsc{Gauss} zu stammen scheinen.

Der zweite Teil der Notiz [6] wurde mit aufgenommen, obwohl er in keinem direkten Zusammenhang mit den Art. 88---90 steht. Man vergleiche Band VIII, S. 128.

Die Formeln der Notiz [15] finden sich bereits in der Disquisitio de elementis ellipticis Palladis (Bd. VI, S. 17---18). Sie wurden hier in eine für die numerische Rechnung geeignetere Form gebracht. Auch das Ergebnis des numerischen Beispiels findet sich dort auf S. 19. Die Abweichung beim Wert von \(C\) rührt daher, dass in der Notiz [15] bei der Herausgabe ein Rechenfehler korrigiert wurde.

Zu den Grenzen der geozentrischen Örter der Himmelskörper (Zodiakus, Notiz [16]---[18]) finden sich im Nachlass auf losen Blättern noch weitere einzelne Herleitungen und Rechnungen. Sie erwiesen sich jedoch als zu fragmentarisch für eine Bearbeitung.

An mehreren Stellen mussten geringfügige Änderungen am Originaltext vorgenommen werden.
\begin{flushright}
\textsc{Brendel}.
\end{flushright}
\begin{center}
\rule{0.3\linewidth}{0.4pt}
\end{center}
